\section{Discussion}
\label{sec:discussion}
\subsection{Predictive transfer and gene-list consistency are distinct}
\label{sec:discussion_tradeoff}
This study tested whether feature-selection strategies that predict CD in an unseen cohort also return consistent gene lists. Under a fixed logistic classifier, TENTACLES V2 had the highest observed mean AUROC (0.723), but the method leading on AUROC changed across D1--D4. The cohort-specific results indicate substantial heterogeneity, and a 0.011 difference between TENTACLES and Elastic Net in four-cohort mean AUROC does not establish that either is generally superior. Patient-stratified bootstrap intervals quantify sampling variability within each test cohort, not uncertainty about deployment to new cohorts.

Gene agreement told a different story. LASSO and TENTACLES produced nearly identical mean LOCO Jaccards (0.211 and 0.210) with markedly different panel sizes; Elastic Net had higher predictive AUROC than LASSO but lower gene Jaccard. Thus neither model performance nor panel size alone determines agreement. Correlated predictors may substitute for one another, while regularisation strength and method-specific selection objectives can change signature sizes. The controlled downstream predictor holds the classifier and test partitions constant, but the upstream tuning procedures, preprocessing and number of genes still vary, so differences cannot be attributed causally to the selector alone.

\subsection{What reproducibility can be inferred?}
\label{sec:discussion_reproducibility}
The six Jaccard comparisons per method were calculated between LOCO fits that share two of their three training cohorts. They therefore quantify how sensitive each selector is to which cohort was excluded, rather than providing four independent discovery replications. Indeed, the separate LASSO and Elastic Net analyses fitted to patient-disjoint cohorts selected thousands of genes and showed substantial size-dependent overlap, illustrating why raw Jaccard is not sufficient to identify stable, selective clinical markers. The fixed top-four analysis partly controls signature length but covers only four selectors and uses method-specific importance rankings; small lists are especially sensitive to ties and individual genes. Native TENTACLES votes have many ties, so an arbitrary top-four vote ranking would not produce a defensible five-method fairness result.

TENTACLES selected six genes in every LOCO fit, including \textit{NOS2} and \textit{CLDN18}, but repeated selection from overlapping training sets is not a clinical validation of any of them. Confirming a candidate biomarker would require a separately collected clinical cohort, consistent control definitions and direct assessment of diagnostic value.

\subsection{Pathway-level patterns are conditional on panel size}
\label{sec:discussion_pathway}
Pathway enrichment sometimes revealed more recurring terms than exact gene-list overlap. In particular, the corrected TENTACLES V2 panels showed recurring interleukin and GPCR-related annotations across D2--D4, while its D1 panel had none passing FDR correction. Elastic Net produced significant terms in all four folds, yet the significant-pathway Jaccard remained moderate (0.273). This does not establish that pathways are universally more reproducible than genes: pathway terms are correlated, the tests use highly unequal gene panels and large panels often have greater statistical power. Stability Selection and StableMate's lack of significant enriched pathways may reflect their 4--20-gene signatures, not the absence of relevant biology. An empty union of significant pathways yields an undefined Jaccard, not evidence of either perfect consistency or zero biological similarity.

\subsection{Cohort composition and implementation limitations}
\label{sec:discussion_limits}
The four analytical cohorts draw on six BioProjects, not six separate holdouts. They vary in CD prevalence, intestinal tissue site, comparison-group composition and assigned count totals. D3 comparison tissue came from colorectal-cancer surgery; tissue site is missing from the supplied D4 metadata. Gene expression may therefore encode non-disease differences that transfer unevenly between cohorts. Descriptive PCA shows cohort structure but cannot distinguish true biological shifts from technical batch effects. We did not apply post-hoc batch correction using the held-out cohorts.

Seven D3 source records had an exclusion flag. Retaining them preserves the published TENTACLES four-cohort collection, but leaves an unperformed QC-exclusion sensitivity analysis. The original RNA-seq reads and clinical diagnoses were not re-audited; only the supplied count tables, metadata alignment, signatures and evaluation outputs were checked.

The V2 output archive independently reproduced the selected nine-voter consensus lists, AUROC, balanced accuracy, and patient-level predictions, and the fixed logistic predictor could be refitted to reproduce V2 probabilities within numerical tolerance. This is not equivalent to independently rerunning all ten author-native classifiers from their original large R objects. The saved D3 and D4 recoveries preserved the fitted models and excluded XGBoost from consensus as prespecified; D4 XGBoost had only four selected cross-validation resamples. Warnings involving native MLP, Boruta and LightGBM candidates were recorded but were not proven harmless. We consequently interpret the nine-voter panel results as an audited adaptation of TENTACLES, not an exact replication of the original publication or proof that every native tuning candidate succeeded.

\subsection{Implications}
\label{sec:discussion_implications}
Transcriptomic benchmark reports should give held-out performance by source cohort alongside the size and overlap of selected signatures, rather than publishing a single mean accuracy or one consensus gene list. Repeating the experiment in independent, tissue-matched cohorts, testing prespecified panel sizes and sensitivity to D3's exclusion flags would better separate biological reproducibility from training-set dependence. Within this four-cohort benchmark, the results support an empirical comparison of candidate panels, but not a clinically validated biomarker claim.

\section{Conclusion}
\label{sec:conclusion}
Across 689 intestinal RNA-seq samples, five selection strategies gave cohort-dependent predictive performance under the same downstream classifier and markedly different gene panels. The highest observed mean AUROC (0.723, TENTACLES V2) did not imply uniformly superior performance or independent biomarker reproducibility. Gene-set overlap and Reactome enrichment also depended on overlapping discovery sets, panel size and source-cohort differences. Cross-cohort prediction, selection consistency and biological validation should therefore be reported and interpreted as distinct outcomes.
